\section{Comparison of the Corrected 2015 Polynomial Re-analysis and the GPR Method}
\label{app:2015-method-comparison}

This appendix compares a historical 95\% CL GPR-based 2015 upper limit with the corrected
2015 re-analysis documented in Appendix~J of the 2016 bump-hunt internal note
\cite{HPS2016BumpInternal}.  The comparison is useful because both analyses use the
same fully unblinded 2015 engineering-run mass spectrum, the same basic prompt
$A^\prime\to e^+e^-$ signal hypothesis, and closely related detector inputs for the
mass resolution and radiative normalization.  At the same time, the comparison should
not be interpreted as a proof that the two statistical procedures are equivalent.  The
corrected 2015 re-analysis is a local windowed polynomial fit, while the present
analysis uses a sideband-trained Gaussian process to predict the blinded mass window
and propagates the resulting correlated background uncertainty through the profiled
likelihood.

\subsection{Historical 2015 and 2016 polynomial analyses}

The published 2015 prompt resonance search used the full 2015 engineering-run data set
passing the analysis selection and scanned the invariant-mass spectrum from roughly
\SIrange{19}{81}{MeV} \cite{HPS2015DarkPhoton,HPS2015InternalNote}.  For each mass
hypothesis, the local model consisted of a Gaussian signal component plus a smooth
background described by a Chebyshev polynomial of the first kind.  The published
optimization selected a fifth-order polynomial below about \SI{39}{MeV} and a
third-order polynomial above that mass, with broad mass windows tied to the detector
resolution.  Upper limits were obtained with a profile-likelihood construction and a
power-constrained prescription: when the unconstrained observed limit became stronger
than the median background-only sensitivity, the quoted signal-yield limit was set to
the median expected limit \cite{HPS2015InternalNote}.  The corresponding coupling
limit used the standard fixed-target relation between the signal yield and the
radiative trident rate, with the 2015 radiative fraction convention
$f_{\rm rad}=0.085$.

The 2016 bump-hunt analysis generalized the polynomial approach.  It used a Gaussian
signal component plus a Legendre-polynomial background in a mass window scaled to
$[-1,1]$, so that the polynomial basis was orthogonal over the fitted interval.  The
polynomial order and window size were selected separately by mass, using 10\% data and
toy studies to check the background-only fit quality, the fitted signal pull, and the
expected limit behavior.  The 2016 note then used a bounded profile-likelihood
\CLs{} construction for the upper limit, with systematic variations propagated by
repeated fits or conservative conversion factors \cite{HPS2016BumpInternal}.  The
``Legendre Polynomial Approach'' curve used in Fig.~\ref{fig:hps2015-legendre-gpr-observed}
is the corrected 2015 re-analysis from that 2016 code path and Appendix~J, not the
original published 2015 contour.

\subsection{Implementation issues in the original 2015 result}

Appendix~J of the 2016 internal note identifies two separate implementation errors in
the original 2015 bump-hunt code \cite{HPS2016BumpInternal}.  First, the local fit
window was scaled to $[-0.25,0.25]$ rather than to the intended $[-1,1]$.  This was not
a harmless convention change: the polynomial background basis was chosen on the
assumption that it would be orthogonal over the fitted coordinate range, so the
incorrect scaling broke the orthogonality argument and invalidated the original
model-selection basis.  Second, the mass resolution was parameterized in the unscaled
mass coordinate but was not consistently transformed into the scaled fit coordinate.
The signal template width entering the fit was therefore inconsistent with the
coordinate system used by the likelihood.

After these corrections, the model selection preferred narrower and more stable
windows for the 2015 re-analysis:
\begin{equation}
M(N,n_\sigma)=
\begin{cases}
O(5),~n_\sigma=10, & m_{A^\prime}<\SI{40}{MeV},\\
O(3),~n_\sigma=8, & m_{A^\prime}\ge \SI{40}{MeV}.
\end{cases}
\end{equation}
The corrected re-analysis found no significant excess and produced the revised
$\eps^2$ limit shown in Fig.~71 of the 2016 internal note.  Figure~72 of that note
compares the original and corrected 2015 contours directly.

This issue should be distinguished from the RooFit binned extended-maximum-likelihood
problem discussed in the 2015 internal note \cite{HPS2015InternalNote}.  That RooFit
problem produced biases in high-occupancy binned fits with several floating
parameters, but the 2015 internal note states that the final results used the
independent vanilla ROOT fitter.  The Appendix~J corrections are therefore a separate
coordinate-scaling and signal-width consistency problem in the 2015 resonance-search
implementation.

\subsection{Present GPR method}

The GPR method is not a polynomial re-fit with a smoother background curve.
For each test mass, the method excludes the blinded region from the training data,
fits the smooth invariant-mass spectrum in log-count space, and predicts the
count-space mean vector and covariance matrix in the blinded mass window.  The local
signal extraction then uses a multi-bin Poisson likelihood in which the GP-predicted
background covariance is represented by correlated Gaussian-constrained nuisance
parameters.  A full-range Gaussian signal template is normalized over the histogram
support and then sliced into the blinded window without re-normalization, so the fitted
signal amplitude remains the total local Gaussian yield associated with the mass
hypothesis. The comparison curve uses a 95\% CL upper limit obtained with the bounded
profile-likelihood \CLs{} construction described in Section~\ref{sec:methodology}; it
is retained as a historical cross-check and is not the v3 90\% CL candidate.

For the observed-limit comparison shown here, the GPR curve corresponds to an earlier
2015 configuration with a widened $2.25\,\sigma_m$ blind and training
exclusion region, a \SI{1}{MeV} scan step, the 7\% radiative-normalization penalty, and
the $\eps^2$ density evaluated with the \SI{1.64}{\sigma_m} convention.  The
validation studies in Section~\ref{sec:toys} motivate the widened blind/extraction
window through leakage, pull, amplitude-recovery, and significance-calibration
diagnostics. This figure is not used to define the v3 confidence-level convention or
the simultaneous candidate.

\begin{table}[p]
\centering
\small
\begin{tabularx}{\linewidth}{P{0.25\linewidth}YY}
\toprule
Feature & Corrected 2015 polynomial re-analysis & Historical GPR comparison \\
\midrule
Data set & Fully unblinded 2015 engineering-run spectrum over the published low-mass
range. & Same 2015 spectrum for the comparison; the v3 scan begins at the published
lower boundary of \SI{19}{MeV} and extends the upper boundary from \SI{81}{MeV} to
\SI{90}{MeV}. \\
Background model & Local mass-window fit with a Legendre-polynomial background in the
corrected 2016-code re-analysis; original 2015 publication used Chebyshev
polynomials. & Full-range log-count GP trained outside the blinded region, with a
count-space blind-window mean and covariance. \\
Signal model & Gaussian resonance template with mass resolution from the 2015 detector
model. & Full-range Gaussian template sliced into the blinded window without
renormalization. \\
Uncertainty treatment & Polynomial coefficients profiled locally; mass-resolution and
radiative-fraction effects propagated by refits, quantiles, or conversion scaling. &
Correlated GP background uncertainty enters directly as Gaussian-constrained nuisance
parameters in the local likelihood; radiative fraction currently enters as a
conversion penalty. \\
Limit construction & Original 2015 publication used a power-constrained 95\% C.L.
limit; the 2016-code re-analysis uses the later bounded \CLs{} machinery. & Bounded
profile-likelihood \CLs{} limit in signal-yield space, converted to $\eps^2$ with the
dataset-specific radiative normalization and local event density. \\
Validation role & Tests whether the corrected historical polynomial implementation
gives a stable 2015 result. & Requires independent toy closure, coverage, and
signal-absorption studies in addition to observed-limit agreement. \\
\bottomrule
\end{tabularx}
\caption{Methodological comparison between the corrected 2015 polynomial
re-analysis and the present GPR analysis.}
\label{tab:2015-polynomial-gpr-methods}
\end{table}

\subsection{Observed-limit comparison}

Figure~\ref{fig:hps2015-legendre-gpr-observed} overlays the observed 95\% upper limit
from the corrected 2015 polynomial re-analysis with the observed 2015 GPR limit.  The
historical curve was digitized from the solid black contour in Fig.~71 of the 2016
internal note, with the digitized points and overlay QA archived together with the
comparison plot.  Over the common \SIrange{20}{81}{MeV} interval, the ratio
\[
R(m) = \frac{\eps^2_{95,\mathrm{GPR}}(m)}
{\eps^2_{95,\mathrm{Legendre}}(m)}
\]
has a median value of 0.99 and a mean value of 1.06 on the integer-MeV comparison
grid.  The central 68\% interval of the ratio is approximately
$0.73<R<1.42$, and all but two grid points lie within a factor of two.  Thus the two
observed limits are similar at the level expected for analyses using the same data
sample, mass resolution scale, radiative normalization, and local event-density
conversion.

\begin{figure}[p]
\centering
\graphicorplaceholder{0.97\linewidth}{final_limit_projection_figs/individual_datasets/hps2015_95cl_observed_vs_internal_note.png}
\caption{Observed 95\% $\eps^2$ upper-limit comparison for the 2015 data set.  The
black curve labeled ``Legendre Polynomial Approach'' is the corrected 2015
re-analysis digitized from Fig.~71 of the 2016 bump-hunt internal note.  The blue curve
labeled ``GPR approach'' is the historical 95\% CL GPR result for the same data set.
The lower panel shows the pointwise GPR/Legendre ratio; over the common
\SIrange{20}{81}{MeV} grid its median is 0.99 and its mean is 1.06. Agreement at this
level checks the overall yield-to-$\eps^2$ normalization and observed sensitivity
scale, but does not by itself test background-model coverage or signal absorption.}
\label{fig:hps2015-legendre-gpr-observed}
\end{figure}

\subsection{Interpretation}

The close observed-limit agreement is reassuring, but its meaning is limited.  An
observed $\eps^2$ limit is driven by the local event density, the radiative-fraction
normalization, the mass-resolution model, and the realized upward or downward
fluctuations in the same data.  Those inputs are largely shared between the two 2015
curves, so agreement in the gross sensitivity scale is expected if neither background
model introduces a large bias.  The agreement therefore supports the statement that the
GPR procedure is operating on the correct physical normalization scale and is not
obviously over- or under-excluding the 2015 data relative to the corrected historical
analysis.

The agreement is not, by itself, a validation of coverage, background-model bias, or
absence of signal absorption. Those questions are addressed by the toy program:
functional-form pseudoexperiments test closure against smooth analytic spectra,
GP-propagated pseudoexperiments test behavior under the GP background model, and
refit-on-toy studies test whether the sideband training can absorb signal-like
structure.  The comparison in Fig.~\ref{fig:hps2015-legendre-gpr-observed} should
therefore be presented as a cross-check of the observed sensitivity scale. The
functional-form closure studies and the outstanding production-matched coverage test
provide the relevant calibration evidence and gate, respectively.

The corrected 2015 re-analysis also changes how the published 2015 result should be
used in this note.  The original published contour should not be used as the primary
target for validating the GPR background model, because Appendix~J shows that the
original coordinate scaling and signal-width treatment were inconsistent with the fit
model assumptions.  The corrected re-analysis is the more appropriate historical
baseline.  At the same time, the corrected analysis did not reveal a significant
2015 excess, so the physics conclusion that the 2015 engineering run did not contain a
discovery-level prompt resonance remains unchanged.  The practical implication is that
comparisons to 2015 should emphasize the corrected polynomial result for methodology
checks, while the published 2015 exclusion remains part of the historical HPS record.
